\documentclass[11pt]{article}
\usepackage{amsmath,amssymb}
\usepackage{graphicx}

\title{Deep Residual Learning for Image Classification}
\author{K. Tanaka \and M. Almeida \and S. Rao}
\date{}

\begin{document}
\maketitle

\begin{abstract}
We revisit residual learning for large-scale image classification. By stacking
identity-skip blocks, gradients propagate cleanly through networks of more than
one hundred layers. Our 152-layer model attains $4.49\%$ top-5 error on
ImageNet validation, narrowing the gap to human-level discrimination on the
1000-way classification task. We further show that residual depth, not width,
is the main driver of accuracy in this regime.
\end{abstract}

\section{Introduction}
Convolutional networks dominate visual recognition, but training them past
twenty layers traditionally degrades accuracy --- a symptom of vanishing rather
than overfitting. The residual reformulation $\mathcal{F}(x) + x$ converts the
optimisation problem into one of learning a perturbation around the identity.
This change is mild in expression but dramatic in practice: depth becomes a
free parameter rather than an obstacle.

\section{Architecture}
Each residual block contains two $3\times3$ convolutions with batch normalisation
and a ReLU non-linearity:
\begin{equation}
y = \sigma\!\left(\mathrm{BN}(W_2 \, \sigma(\mathrm{BN}(W_1 x))) + x\right).
\end{equation}
We downsample with strided convolutions rather than pooling, and project the
shortcut with a $1\times1$ convolution whenever the channel count changes.
Bottleneck variants compress the inner representation through a $1\times1$
$\to 3\times3 \to 1\times1$ stack, trading a slim middle dimension for a wider
input/output.

\section{Training}
We train for 90 epochs of SGD with momentum 0.9, weight decay $10^{-4}$, and a
cosine learning-rate schedule beginning at 0.1. Inputs are resized so the
shorter side is sampled uniformly from $[256, 480]$, then a $224\times224$
crop is taken with random horizontal flips. Mixed precision halves wall-clock
time without measurable accuracy loss.

\section{Results}
\begin{tabular}{lcc}
\hline
Model & Layers & Top-5 Error \\
\hline
Plain & 34  & $7.40\%$ \\
ResNet & 34  & $5.71\%$ \\
ResNet & 152 & $4.49\%$ \\
\hline
\end{tabular}

The plain 34-layer network underperforms its 18-layer counterpart, reproducing
the degradation phenomenon. The residual variant inverts the trend: deeper is
strictly better up to the point where memory becomes the binding constraint.

\section{Discussion}
Skip connections are sometimes characterised as ensemble averaging over
sub-paths. We instead view them as a conditioning device that keeps the
Jacobian's spectral radius near unity throughout training. Either lens is
consistent with the observation that residual depth scales gracefully where
plain depth does not.

\bibliographystyle{plain}
\begin{thebibliography}{9}
\bibitem{he2016deep} K. He et al. \emph{Deep Residual Learning for Image Recognition.} CVPR, 2016.
\bibitem{ioffe2015batch} S. Ioffe, C. Szegedy. \emph{Batch Normalization.} ICML, 2015.
\end{thebibliography}

\end{document}
